\scanpage{50}
\mspage{4}
It follows from the definition that the isomorphism $\hat F:\underline{1}'\xrightarrow{\sim}F\underline{1}$, if it exists, is unique.

A \uline{$\otimes$-AC functor} is \unsure{an|a} $\otimes$-associative and commutative functor. A \uline{$\otimes$-ACU functor} is a $\otimes$-associative, commutative and unifer functor.

Let $(F,\breve F)$ and $(G,\breve G)$ be $\otimes$-functors from a $\otimes$-category $\underline{C}$ to a $\otimes$-category $\underline{C}'$. A \uline{$\otimes$-morphism} from the $\otimes$-functor $(F,\breve F)$ to the $\otimes$-functor $(G,\breve G)$ is a morphism of functors $\lambda:F\to G$ such that the following diagram commutes
\diagram{\begin{tikzcd}[ampersand replacement=\&]
FX\otimes FY \arrow[r,"{\breve F}"] \arrow[d,"{\lambda_X\otimes\lambda_Y}"'] \& F(X\otimes Y) \arrow[d,"{\lambda_{X\otimes Y}}"]\\
GX\otimes GY \arrow[r,"{\breve G}"] \& G(X\otimes Y)
\end{tikzcd}\qquad $X,Y\in\operatorname{Ob}\underline{C}$}

Chapter $\overline{\underline{\mathrm{II}}}$ is a study of Gr-categories and Pic-categories. A \uline{Gr-category} is a $\otimes$-AU category, the objects of which are all invertible, and the base category a groupoïd (i.e all arrows are isomorphisms). Thus a Gr-category is like a group. We obtain from this definition that if $\underline{P}$ is a Gr-category, the set $\Pi_0(\underline{P})$ of the classes up to ismomorphism of objects of $\underline{P}$, together with the operation induced by the law $\otimes$ of $\underline{P}$, is a group; the group $\operatorname{Aut}(\underline{1})=\Pi_1(\underline{P})$ is a commutative group; and for all $X\in\operatorname{Ob}\underline{P}$
\[
\begin{aligned}
\gamma_X:\ u\longmapsto u\otimes\mathrm{id}_X&=\operatorname{Aut}(\underline{1})\xrightarrow{\sim}\operatorname{Aut}(X)\\
\delta_X:\ u\longmapsto\mathrm{id}_X\otimes u&=\operatorname{Aut}(\underline{1})\xrightarrow{\sim}\operatorname{Aut}(X)
\end{aligned}
\]
We attribute thus to a Gr-category $\underline{P}$ two groups $\Pi_0(\underline{P})$ and $\Pi_1(\underline{P})$ where $\Pi_1(\underline{P})$ is commutative. Furthermore we can define an action of $\Pi_0(\underline{P})$ on $\Pi_1(\underline{P})$ by the formula
\[
su=\delta_X^{-1}\gamma_X(u)
\]
for $s\in\Pi_0(\underline{P})$ represented by $X$ and $u\in\Pi_1(\underline{P})$. The commutative group $\Pi_1(\underline{P})$ together with this action is a left $\Pi_0(\underline{P})$-module.

Let $M$ be a group, $N$ a left $M$-module. A \uline{preepinglages} of
% archive sheet number: 49
